\section{Состав и порядок испытаний}

\subsection{Технические средства}

Испытания должны проводиться на вычислительном устройстве, обеспечивающем
запуск 64-разрядных приложений, выполнение командной строки операционной
системы Linux и хранение выходных файлов, формируемых в процессе генерации.

При проведении испытаний, связанных с большими значениями масштаба генерации,
технические средства должны обеспечивать наличие вычислительных ресурсов,
достаточных для формирования и сохранения результирующих наборов данных.

\subsection{Программные средства}

При проведении испытаний должны использоваться следующие программные средства:
\begin{enumerate}
    \item исполняемый файл программы \texttt{schengen\_main};
    \item штатные средства командной строки операционной системы;
    \item файловая система или доступное оператору хранилище для записи
    результатов генерации.
\end{enumerate}

\subsection{Порядок проведения испытаний}

Испытания рекомендуется проводить последовательно, переходя от проверки общих
условий запуска программы к проверке ее основных режимов работы, параметров
управления и диагностических сообщений.

Испытания следует проводить в следующем порядке:
\begin{enumerate}
    \item проверка запуска программы и вывода справочной информации;
    \item проверка вывода списка таблиц и списка доступных форматов;
    \item проверка соответствия генерации эталонным данным набора TPC-H;
    \item проверка генерации полного набора таблиц на малом масштабе;
    \item проверка генерации выбранного подмножества таблиц;
    \item проверка обработки параметров записи и пути вывода;
    \item проверка устойчивой работы программы при больших значениях масштаба
    генерации;
    \item проверка диагностирования ошибочных параметров;
    \item проверка соответствия программной документации фактическому поведению
    программы.
\end{enumerate}
